% Options for packages loaded elsewhere
\PassOptionsToPackage{unicode}{hyperref}
\PassOptionsToPackage{hyphens}{url}
\documentclass[
  american,
]{article}
\usepackage{xcolor}
\usepackage{amsmath,amssymb}
\setcounter{secnumdepth}{5}
\usepackage{iftex}
\ifPDFTeX
  \usepackage[T1]{fontenc}
  \usepackage[utf8]{inputenc}
  \usepackage{textcomp} % provide euro and other symbols
\else % if luatex or xetex
  \usepackage{unicode-math} % this also loads fontspec
  \defaultfontfeatures{Scale=MatchLowercase}
  \defaultfontfeatures[\rmfamily]{Ligatures=TeX,Scale=1}
\fi
\usepackage{lmodern}
\ifPDFTeX\else
  % xetex/luatex font selection
  \setmainfont[]{DejaVu Sans}
  \setsansfont[]{DejaVu Sans}
  \setmonofont[]{DejaVu Sans Mono}
\fi
% Use upquote if available, for straight quotes in verbatim environments
\IfFileExists{upquote.sty}{\usepackage{upquote}}{}
\IfFileExists{microtype.sty}{% use microtype if available
  \usepackage[]{microtype}
  \UseMicrotypeSet[protrusion]{basicmath} % disable protrusion for tt fonts
}{}
\makeatletter
\@ifundefined{KOMAClassName}{% if non-KOMA class
  \IfFileExists{parskip.sty}{%
    \usepackage{parskip}
  }{% else
    \setlength{\parindent}{0pt}
    \setlength{\parskip}{6pt plus 2pt minus 1pt}}
}{% if KOMA class
  \KOMAoptions{parskip=half}}
\makeatother
\usepackage{color}
\usepackage{fancyvrb}
\newcommand{\VerbBar}{|}
\newcommand{\VERB}{\Verb[commandchars=\\\{\}]}
\DefineVerbatimEnvironment{Highlighting}{Verbatim}{commandchars=\\\{\}}
% Add ',fontsize=\small' for more characters per line
\newenvironment{Shaded}{}{}
\newcommand{\AlertTok}[1]{\textcolor[rgb]{1.00,0.00,0.00}{\textbf{#1}}}
\newcommand{\AnnotationTok}[1]{\textcolor[rgb]{0.38,0.63,0.69}{\textbf{\textit{#1}}}}
\newcommand{\AttributeTok}[1]{\textcolor[rgb]{0.49,0.56,0.16}{#1}}
\newcommand{\BaseNTok}[1]{\textcolor[rgb]{0.25,0.63,0.44}{#1}}
\newcommand{\BuiltInTok}[1]{\textcolor[rgb]{0.00,0.50,0.00}{#1}}
\newcommand{\CharTok}[1]{\textcolor[rgb]{0.25,0.44,0.63}{#1}}
\newcommand{\CommentTok}[1]{\textcolor[rgb]{0.38,0.63,0.69}{\textit{#1}}}
\newcommand{\CommentVarTok}[1]{\textcolor[rgb]{0.38,0.63,0.69}{\textbf{\textit{#1}}}}
\newcommand{\ConstantTok}[1]{\textcolor[rgb]{0.53,0.00,0.00}{#1}}
\newcommand{\ControlFlowTok}[1]{\textcolor[rgb]{0.00,0.44,0.13}{\textbf{#1}}}
\newcommand{\DataTypeTok}[1]{\textcolor[rgb]{0.56,0.13,0.00}{#1}}
\newcommand{\DecValTok}[1]{\textcolor[rgb]{0.25,0.63,0.44}{#1}}
\newcommand{\DocumentationTok}[1]{\textcolor[rgb]{0.73,0.13,0.13}{\textit{#1}}}
\newcommand{\ErrorTok}[1]{\textcolor[rgb]{1.00,0.00,0.00}{\textbf{#1}}}
\newcommand{\ExtensionTok}[1]{#1}
\newcommand{\FloatTok}[1]{\textcolor[rgb]{0.25,0.63,0.44}{#1}}
\newcommand{\FunctionTok}[1]{\textcolor[rgb]{0.02,0.16,0.49}{#1}}
\newcommand{\ImportTok}[1]{\textcolor[rgb]{0.00,0.50,0.00}{\textbf{#1}}}
\newcommand{\InformationTok}[1]{\textcolor[rgb]{0.38,0.63,0.69}{\textbf{\textit{#1}}}}
\newcommand{\KeywordTok}[1]{\textcolor[rgb]{0.00,0.44,0.13}{\textbf{#1}}}
\newcommand{\NormalTok}[1]{#1}
\newcommand{\OperatorTok}[1]{\textcolor[rgb]{0.40,0.40,0.40}{#1}}
\newcommand{\OtherTok}[1]{\textcolor[rgb]{0.00,0.44,0.13}{#1}}
\newcommand{\PreprocessorTok}[1]{\textcolor[rgb]{0.74,0.48,0.00}{#1}}
\newcommand{\RegionMarkerTok}[1]{#1}
\newcommand{\SpecialCharTok}[1]{\textcolor[rgb]{0.25,0.44,0.63}{#1}}
\newcommand{\SpecialStringTok}[1]{\textcolor[rgb]{0.73,0.40,0.53}{#1}}
\newcommand{\StringTok}[1]{\textcolor[rgb]{0.25,0.44,0.63}{#1}}
\newcommand{\VariableTok}[1]{\textcolor[rgb]{0.10,0.09,0.49}{#1}}
\newcommand{\VerbatimStringTok}[1]{\textcolor[rgb]{0.25,0.44,0.63}{#1}}
\newcommand{\WarningTok}[1]{\textcolor[rgb]{0.38,0.63,0.69}{\textbf{\textit{#1}}}}
\usepackage{longtable,booktabs,array}
\usepackage{calc} % for calculating minipage widths
% Correct order of tables after \paragraph or \subparagraph
\usepackage{etoolbox}
\makeatletter
\patchcmd\longtable{\par}{\if@noskipsec\mbox{}\fi\par}{}{}
\makeatother
% Allow footnotes in longtable head/foot
\IfFileExists{footnotehyper.sty}{\usepackage{footnotehyper}}{\usepackage{footnote}}
\makesavenoteenv{longtable}
\usepackage{graphicx}
\makeatletter
\newsavebox\pandoc@box
\newcommand*\pandocbounded[1]{% scales image to fit in text height/width
  \sbox\pandoc@box{#1}%
  \Gscale@div\@tempa{\textheight}{\dimexpr\ht\pandoc@box+\dp\pandoc@box\relax}%
  \Gscale@div\@tempb{\linewidth}{\wd\pandoc@box}%
  \ifdim\@tempb\p@<\@tempa\p@\let\@tempa\@tempb\fi% select the smaller of both
  \ifdim\@tempa\p@<\p@\scalebox{\@tempa}{\usebox\pandoc@box}%
  \else\usebox{\pandoc@box}%
  \fi%
}
% Set default figure placement to htbp
\def\fps@figure{htbp}
\makeatother
% definitions for citeproc citations
\NewDocumentCommand\citeproctext{}{}
\NewDocumentCommand\citeproc{mm}{%
  \begingroup\def\citeproctext{#2}\cite{#1}\endgroup}
\makeatletter
 % allow citations to break across lines
 \let\@cite@ofmt\@firstofone
 % avoid brackets around text for \cite:
 \def\@biblabel#1{}
 \def\@cite#1#2{{#1\if@tempswa , #2\fi}}
\makeatother
\newlength{\cslhangindent}
\setlength{\cslhangindent}{1.5em}
\newlength{\csllabelwidth}
\setlength{\csllabelwidth}{3em}
\newenvironment{CSLReferences}[2] % #1 hanging-indent, #2 entry-spacing
 {\begin{list}{}{%
  \setlength{\itemindent}{0pt}
  \setlength{\leftmargin}{0pt}
  \setlength{\parsep}{0pt}
  % turn on hanging indent if param 1 is 1
  \ifodd #1
   \setlength{\leftmargin}{\cslhangindent}
   \setlength{\itemindent}{-1\cslhangindent}
  \fi
  % set entry spacing
  \setlength{\itemsep}{#2\baselineskip}}}
 {\end{list}}
\usepackage{calc}
\newcommand{\CSLBlock}[1]{\hfill\break\parbox[t]{\linewidth}{\strut\ignorespaces#1\strut}}
\newcommand{\CSLLeftMargin}[1]{\parbox[t]{\csllabelwidth}{\strut#1\strut}}
\newcommand{\CSLRightInline}[1]{\parbox[t]{\linewidth - \csllabelwidth}{\strut#1\strut}}
\newcommand{\CSLIndent}[1]{\hspace{\cslhangindent}#1}
\ifLuaTeX
\usepackage[bidi=basic,shorthands=off,]{babel}
\else
\usepackage[bidi=default,shorthands=off,]{babel}
\fi
\ifPDFTeX
\else
\babelfont{rm}[]{DejaVu Sans}
\fi
\ifLuaTeX
  \usepackage{selnolig} % disable illegal ligatures
\fi
\setlength{\emergencystretch}{3em} % prevent overfull lines
\providecommand{\tightlist}{%
  \setlength{\itemsep}{0pt}\setlength{\parskip}{0pt}}
\usepackage{fontspec}
\usepackage{xcolor}
\usepackage{tcolorbox}
\usepackage{booktabs}
\usepackage{longtable}
\usepackage{array}
\usepackage{caption}
\usepackage{float}
\usepackage{fancyhdr}
\usepackage{tocloft}
\usepackage{geometry}
\geometry{margin=1in}
\definecolor{mtcred}{HTML}{E30613}
\definecolor{lightblue}{HTML}{EEF7FF}
\definecolor{lightred}{HTML}{FFF1F3}
\definecolor{lightgraybox}{HTML}{F6F7F9}
\definecolor{darktext}{HTML}{222222}
\AtBeginDocument{\hypersetup{colorlinks=true, linkcolor=mtcred, urlcolor=mtcred, citecolor=mtcred}}
\pagestyle{fancy}
\fancyhf{}
\lhead{n-step SAC final project}
\rhead{Radmir I.}
\cfoot{\thepage}
\setlength{\headheight}{14pt}
\newtcolorbox{keybox}{
  colback=lightblue,
  colframe=mtcred!70!black,
  boxrule=0.6pt,
  arc=3mm,
  left=2mm,
  right=2mm,
  top=1mm,
  bottom=1mm
}
\newtcolorbox{notegray}{
  colback=lightgraybox,
  colframe=black!15,
  boxrule=0.4pt,
  arc=3mm,
  left=2mm,
  right=2mm,
  top=1mm,
  bottom=1mm
}
\newtcolorbox{warningbox}{
  colback=lightred,
  colframe=mtcred,
  boxrule=0.6pt,
  arc=3mm,
  left=2mm,
  right=2mm,
  top=1mm,
  bottom=1mm
}
\captionsetup{font=small, labelfont=bf}

\setlength{\cftsecnumwidth}{3em}
\setlength{\cftsubsecnumwidth}{3.5em}
\usepackage{bookmark}
\IfFileExists{xurl.sty}{\usepackage{xurl}}{} % add URL line breaks if available
\urlstyle{same}
\hypersetup{
  pdftitle={Stress-Testing n-step Soft Actor-Critic under Limited Compute},
  pdfauthor={Radmir I.},
  pdflang={en-US},
  hidelinks,
  pdfcreator={LaTeX via pandoc}}

\title{Stress-Testing n-step Soft Actor-Critic under Limited Compute}
\usepackage{etoolbox}
\makeatletter
\providecommand{\subtitle}[1]{% add subtitle to \maketitle
  \apptocmd{\@title}{\par {\large #1 \par}}{}{}
}
\makeatother
\subtitle{Final RL Project Report, extended version}
\author{Radmir I.}
\date{2026}

\begin{document}
\maketitle

{
\setcounter{tocdepth}{3}
\tableofcontents
}
\begin{keybox}
\textbf{Short version.} This project tests one central mechanism behind recent n-step SAC work: whether replacing the standard 1-step critic target with an n-step target helps Soft Actor-Critic under a limited compute budget. On HalfCheetah-v5, the 3-step variant improved learning compared with the 1-step baseline, while the 5-step variant failed under the same hyperparameters. The main conclusion is deliberately narrow: the n-step horizon behaves like a stability parameter, not like a simple bigger-is-better knob.
\end{keybox}

\section{Abstract}\label{abstract}

This project studies a compact n-step modification of Soft Actor-Critic
(SAC) under a limited compute budget. The main question is whether
replacing the standard 1-step critic target with n-step targets improves
sample efficiency in continuous-control tasks, and where the horizon
becomes unstable. I implement a configurable n-step buffer before the
replay buffer and compare SAC with \(n=1\), \(n=3\), and \(n=5\). The
experiments include Pendulum-v1 as a sanity check and HalfCheetah-v5 as
the main benchmark. On HalfCheetah-v5, \(n=3\) reaches the best final
return and reaches high-return behavior earlier than the baseline, while
\(n=5\) fails to learn a strong running policy under the same
hyperparameters. The result supports the idea that n-step credit
assignment can help SAC, but only within a stability-sensitive range.

\section{Introduction}\label{introduction}

Soft Actor-Critic is a strong off-policy actor-critic algorithm for
continuous-control tasks. It combines replay-based off-policy updates
with a stochastic actor and a maximum-entropy objective, encouraging the
policy to maximize both return and entropy (Haarnoja, Zhou, Abbeel, et
al. 2018; Haarnoja, Zhou, Hartikainen, et al. 2018). This makes SAC a
natural baseline for continuous-control experiments: it is
sample-efficient, relatively stable, and already hard to improve without
careful changes.

In standard SAC, the critic is usually trained with a 1-step
temporal-difference target. This target is simple and stable, but reward
information is propagated backward only one step at a time. N-step
returns propose a more aggressive alternative: accumulate several
rewards before bootstrapping. This can improve credit assignment because
useful reward information can move backward faster (Sutton and Barto
2018).

However, SAC is off-policy. Transitions sampled from the replay buffer
may have been collected by older policies. A longer n-step target
depends on a longer replay fragment and may therefore introduce more
off-policy bias and target variance. This is the core reason why n-step
SAC is non-trivial and why recent work has revisited the topic (Łyskawa,
Lewandowski, and Wawrzyński 2025; Tian, Celik, and Neumann 2025).

The project is intentionally not a state-of-the-art claim. It is a
bounded reproduction and ablation study: I isolate one mechanism - the
critic target horizon - and test whether the expected trade-off appears
in a compact, reproducible setup.

\section{Research question and
contribution}\label{research-question-and-contribution}

The main research question is:

\begin{keybox}
Do n-step returns improve SAC under limited compute, and how sensitive is the result to the chosen horizon?
\end{keybox}

The project compares three variants:

\begin{longtable}[]{@{}lll@{}}
\toprule\noalign{}
Method & Description & Expected behavior \\
\midrule\noalign{}
\endhead
\bottomrule\noalign{}
\endlastfoot
SAC \(n=1\) & standard 1-step SAC baseline & stable baseline \\
SAC \(n=3\) & moderate 3-step target & may improve sample efficiency \\
SAC \(n=5\) & longer 5-step target & potentially higher risk \\
\end{longtable}

The contribution is small but concrete:

\begin{enumerate}
\def\labelenumi{\arabic{enumi}.}
\tightlist
\item
  a configurable n-step transition buffer integrated into a compact SAC
  implementation;
\item
  a controlled \(n=1/3/5\) ablation on Pendulum-v1 and HalfCheetah-v5;
\item
  committed raw logs, plots, summary tables, and reproducibility notes;
\item
  a failure analysis of the \(n=5\) run under fixed hyperparameters.
\end{enumerate}

\section{Related work}\label{related-work}

\subsection{Soft Actor-Critic}\label{soft-actor-critic}

SAC is an off-policy maximum-entropy actor-critic algorithm. The actor
is optimized not only for expected return but also for entropy, which
encourages exploration and can improve robustness. The practical SAC
formulation also uses twin Q-critics, target networks, replay buffer
sampling, and automatic entropy temperature tuning (Haarnoja, Zhou,
Abbeel, et al. 2018; Haarnoja, Zhou, Hartikainen, et al. 2018).

SAC is a good baseline for this project because it is already a strong
continuous-control method. If a small n-step modification helps or
breaks SAC, the effect is meaningful enough to analyze.

\subsection{N-step returns and credit
assignment}\label{n-step-returns-and-credit-assignment}

N-step returns are a classical temporal-difference learning idea.
Instead of bootstrapping after a single reward, an n-step target
accumulates several rewards before bootstrapping:

\[
G_t^{(n)} = r_t + \gamma r_{t+1} + \dots + \gamma^{n-1} r_{t+n-1} + \gamma^n V(s_{t+n}).
\]

This can propagate delayed rewards faster than a 1-step target. At the
same time, the target contains more sampled rewards and changes the
bias-variance trade-off (Sutton and Barto 2018; Sharma et al. 2017).

\subsection{Off-policy multi-step
correction}\label{off-policy-multi-step-correction}

The off-policy setting is the main difficulty. In replay-based learning,
longer trajectory fragments may have been collected by a behavior policy
different from the current target policy. Several multi-step methods
explicitly address this issue through correction or trace truncation.
Retrace(\(\lambda\)) is designed as a safe and efficient off-policy
return-based method (Munos et al. 2016). IMPALA introduces V-trace to
stabilize learning under decoupled acting and learning (Espeholt et al.
2018). These works are not implemented here, but they motivate why naive
off-policy multi-step targets should be treated carefully.

A related warning comes from the ``deadly triad'' of function
approximation, bootstrapping, and off-policy learning (Hasselt et al.
2018). Naive n-step SAC touches all three: it uses neural critics,
bootstrapped targets, and replayed off-policy transitions.

\subsection{SACn and T-SAC}\label{sacn-and-t-sac}

SACn directly studies how to combine Soft Actor-Critic with n-step
returns. Its motivation is close to this project: n-step returns can
speed up learning, but the usual combination with SAC introduces bias
due to changing action distributions, and importance sampling may become
numerically unstable (Łyskawa, Lewandowski, and Wawrzyński 2025).

T-SAC / Chunking the Critic takes a more ambitious direction. It uses a
sequence-conditioned critic with trajectory chunks and aggregated n-step
targets, making temporal context part of critic learning rather than
only changing a scalar target (Tian, Celik, and Neumann 2025).

This project does not reproduce SACn or T-SAC. Instead, it isolates a
smaller shared mechanism: changing the critic target horizon. This makes
the experiment feasible on limited compute and easier to interpret.

\section{Method}\label{method}

\subsection{Standard SAC components}\label{standard-sac-components}

The implementation keeps the usual SAC structure:

\begin{itemize}
\tightlist
\item
  squashed Gaussian actor;
\item
  twin Q-critics;
\item
  target critic networks;
\item
  replay buffer;
\item
  automatic entropy tuning;
\item
  deterministic evaluation checkpoints.
\end{itemize}

The critic update uses the minimum of the two target critics and
subtracts the entropy term:

\[
V_{\text{target}}(s') =
\min_i Q_i^{\text{target}}(s', a') - \alpha \log \pi(a'|s').
\]

For the standard 1-step version, the target is:

\[
y_t = r_t + \gamma (1-d_t) V_{\text{target}}(s_{t+1}).
\]

\subsection{N-step transition buffer}\label{n-step-transition-buffer}

The only algorithmic component varied in the main ablation is the
transition aggregation before the replay buffer. The n-step buffer
stores a short queue of recent transitions. When enough transitions are
available, it builds an aggregated transition:

\begin{Shaded}
\begin{Highlighting}[]
\NormalTok{state      = s\_t}
\NormalTok{action     = a\_t}
\NormalTok{reward\_n   = r\_t + gamma r\_\{t+1\} + ... + gamma\^{}(k{-}1) r\_\{t+k{-}1\}}
\NormalTok{next\_state = s\_\{t+k\}}
\NormalTok{discount\_n = gamma\^{}k}
\NormalTok{done       = terminal flag inside the n{-}step window}
\end{Highlighting}
\end{Shaded}

Here \(k \leq n\). If the episode ends before the full window is
collected, the window is cut early. For \(n=1\), this construction
reduces to standard SAC.

The critic target becomes:

\[
y_t^{(n)} = R_t^{(n)} + (1-d_t) \gamma^k V_{\text{target}}(s_{t+k}).
\]

The replay buffer therefore samples the same kind of objects for all
runs, but the reward and discount fields encode different target
horizons.

\subsection{What is fixed in the
ablation}\label{what-is-fixed-in-the-ablation}

The comparison keeps the same architecture and hyperparameters for
\(n=1\), \(n=3\), and \(n=5\). This choice is intentional. It makes the
experiment a clean sensitivity check of the target horizon.

\begin{notegray}
This also means that the comparison is not fully tuned for every horizon. A separate hyperparameter search might improve the $n=5$ run. The current result should therefore be read as a fixed-configuration stability test, not as proof that $n=5$ can never work.
\end{notegray}

\section{Experimental setup}\label{experimental-setup}

\subsection{Environments}\label{environments}

Two environments are used.

\begin{longtable}[]{@{}llrr@{}}
\toprule\noalign{}
Environment & Role & Steps & Seed \\
\midrule\noalign{}
\endhead
\bottomrule\noalign{}
\endlastfoot
Pendulum-v1 & sanity check & 30,000 & 0 \\
HalfCheetah-v5 & main benchmark & 100,000 & 0 \\
\end{longtable}

Pendulum-v1 is a simple classic-control continuous-action task with
action space \texttt{Box(-2.0,\ 2.0,\ (1,),\ float32)} and observation
shape \texttt{(3,)} (Farama Foundation 2026b). HalfCheetah-v5 is a
MuJoCo continuous-control locomotion task with action space
\texttt{Box(-1.0,\ 1.0,\ (6,),\ float32)} and observation shape
\texttt{(17,)}; the goal is to make the 2D robot run forward as fast as
possible (Farama Foundation 2026a).

\subsection{Hyperparameters}\label{hyperparameters}

\begin{longtable}[]{@{}lr@{}}
\toprule\noalign{}
Parameter & Value \\
\midrule\noalign{}
\endhead
\bottomrule\noalign{}
\endlastfoot
Discount factor \texttt{gamma} & 0.99 \\
Target update coefficient \texttt{tau} & 0.005 \\
Initial entropy coefficient \texttt{alpha} & 0.2 \\
Automatic entropy tuning & yes \\
Actor learning rate & 3e-4 \\
Critic learning rate & 3e-4 \\
Alpha learning rate & 3e-4 \\
Batch size & 256 \\
Replay buffer size & 1,000,000 \\
Hidden sizes & {[}256, 256{]} \\
Start steps & 5,000 \\
Update after & 5,000 \\
Evaluation interval & 5,000 steps \\
Evaluation episodes & 5 \\
Compared horizons & \(n=1,3,5\) \\
\end{longtable}

All completed runs use deterministic evaluation with five episodes per
checkpoint. Raw CSV logs and plots are committed to the repository.

\subsection{Hardware}\label{hardware}

The experiments were run on a Ryzen 5700 laptop through WSL Ubuntu using
CPU training. This is intentionally a limited-compute setup. The project
is designed as a reproducible student experiment rather than a
large-scale benchmark.

\section{Results}\label{results}

\subsection{Pendulum-v1 sanity check}\label{pendulum-v1-sanity-check}

On Pendulum-v1, all three variants train successfully and reach similar
final quality. The main role of this experiment is to verify that the
implementation works and that the n-step target construction does not
immediately break SAC on a simple environment.

\begin{figure}
\centering
\includegraphics[width=0.9\linewidth,height=\textheight,keepaspectratio,alt={Pendulum-v1 learning curves}]{results/plots/pendulum_learning_curves.png}
\caption{Pendulum-v1 learning curves}\label{fig:pendulum}
\end{figure}

The curves suggest that \(n=3\) reaches a good return quickly, but
Pendulum-v1 is too simple to support a strong claim about the best
horizon. This experiment should be treated as a sanity check, not as the
main evidence.

\subsection{HalfCheetah-v5 main
ablation}\label{halfcheetah-v5-main-ablation}

The main result is observed on HalfCheetah-v5.

\begin{figure}
\centering
\includegraphics[width=0.9\linewidth,height=\textheight,keepaspectratio,alt={HalfCheetah-v5 learning curves}]{results/plots/halfcheetah_learning_curves.png}
\caption{HalfCheetah-v5 learning curves}\label{fig:halfcheetah}
\end{figure}

\begin{longtable}[]{@{}lrr@{}}
\toprule\noalign{}
Method & Final eval return & Best eval return \\
\midrule\noalign{}
\endhead
\bottomrule\noalign{}
\endlastfoot
SAC \(n=1\) & 4354.49 & 4354.49 \\
SAC \(n=3\) & 4578.91 & 4578.91 \\
SAC \(n=5\) & 764.49 & 833.37 \\
\end{longtable}

\begin{keybox}
\textbf{Main observation.} In this single-seed run, $n=3$ gives the best final return and reaches strong performance earlier than the 1-step baseline. The $n=5$ variant fails to learn a strong running policy under the same settings.
\end{keybox}

\subsection{Sample-efficiency markers}\label{sample-efficiency-markers}

Final return is not the only useful metric. The HalfCheetah learning
curve also shows a difference in how quickly the variants reach strong
behavior.

\begin{longtable}[]{@{}lrr@{}}
\toprule\noalign{}
Method & Steps to return \(\geq 3000\) & Steps to return
\(\geq 4000\) \\
\midrule\noalign{}
\endhead
\bottomrule\noalign{}
\endlastfoot
SAC \(n=1\) & 65,000 & 100,000 \\
SAC \(n=3\) & 45,000 & 70,000 \\
SAC \(n=5\) & not reached & not reached \\
\end{longtable}

These markers support the visual reading of the curve: \(n=3\) learns
faster in this run, while \(n=5\) never reaches high-return behavior.

\section{Discussion}\label{discussion}

The result supports the expected trade-off. A moderate n-step horizon
can improve sample efficiency because reward information is propagated
backward faster. This likely explains why \(n=3\) improves over the
1-step baseline on HalfCheetah-v5.

At the same time, a longer naive target can increase off-policy bias and
target variance. This likely explains why \(n=5\) fails under the same
hyperparameters. The most useful result is not only that \(n=3\) obtains
the best final return. The more informative result is the failure of
\(n=5\): it shows that increasing the target horizon can change the
stability of off-policy SAC under fixed hyperparameters.

\begin{warningbox}
The safest conclusion is narrow: in this limited-compute single-seed ablation, 3-step SAC improved HalfCheetah learning compared with the 1-step baseline, while 5-step SAC failed under the same hyperparameters. The project does not claim universal superiority of $n=3$.
\end{warningbox}

This failure mode matches the motivation behind recent n-step SAC work.
N-step targets can help credit assignment, but the horizon interacts
with off-policy bias, target variance, critic stability, and entropy
dynamics.

\section{Threats to validity}\label{threats-to-validity}

\subsection{Statistical validity}\label{statistical-validity}

The main limitation is the use of a single seed. The current experiment
is useful as bounded reproduction and ablation evidence, but it is not
enough to make a statistical claim that \(n=3\) is generally better than
\(n=1\), or that \(n=5\) generally fails. A stronger study should use at
least 3-5 seeds and report confidence intervals.

\subsection{Algorithmic validity}\label{algorithmic-validity}

The implementation uses naive n-step targets without off-policy
correction. This is intentional because the project tests the direct
effect of changing the critic target horizon. However, it also means
that the failure of \(n=5\) should not be interpreted as a failure of
all n-step SAC methods. SACn-style correction or sequence-aware critics
could stabilize longer horizons.

\subsection{Hyperparameter fairness}\label{hyperparameter-fairness}

All horizons use the same hyperparameters. This makes the ablation clean
because only the target horizon changes. At the same time, it may
disadvantage \(n=5\). A longer horizon may require different learning
rates, target smoothing, entropy settings, or replay settings.

\subsection{Environment coverage}\label{environment-coverage}

The project uses Pendulum-v1 and HalfCheetah-v5. Pendulum-v1 is mainly a
sanity check; it is too simple to support a strong claim about the best
horizon. HalfCheetah-v5 is the main benchmark, but one MuJoCo task is
not enough for a general conclusion. Future work should add Hopper-v5,
Walker2d-v5, and possibly Ant-v5.

\subsection{Compute budget}\label{compute-budget}

All experiments were run under laptop CPU compute. This shaped the
design: short training budgets, single seed, and compact architecture.
The project should be read as a small reproducible ablation study rather
than a paper-scale benchmark.

\section{Reproducibility checklist}\label{reproducibility-checklist}

\begin{longtable}[]{@{}ll@{}}
\toprule\noalign{}
Item & Status \\
\midrule\noalign{}
\endhead
\bottomrule\noalign{}
\endlastfoot
Code committed & yes \\
Configs committed & yes \\
Raw CSV logs committed & yes \\
Plots committed & yes \\
Summary table committed & yes \\
Random seed recorded & yes, seed 0 \\
Hardware documented & yes \\
Limitations documented & yes \\
Full multi-seed statistics & no \\
\end{longtable}

Key repository artifacts:

\begin{itemize}
\tightlist
\item
  implementation: \texttt{src/rl/nstep\_buffer.py},
  \texttt{src/rl/sac.py}, \texttt{src/train.py};
\item
  configs: \texttt{configs/pendulum\_sac\_n*.yaml},
  \texttt{configs/halfcheetah\_sac\_n*.yaml};
\item
  raw logs: \texttt{results/raw/*.csv};
\item
  plots: \texttt{results/plots/*.png};
\item
  summary: \texttt{results/halfcheetah\_summary.csv};
\item
  slides and notes: \texttt{slides/}, \texttt{docs/}.
\end{itemize}

\section{Future work}\label{future-work}

The most useful extensions are:

\begin{enumerate}
\def\labelenumi{\arabic{enumi}.}
\tightlist
\item
  run 3-5 seeds;
\item
  add more MuJoCo environments such as Hopper-v5, Walker2d-v5, and
  Ant-v5;
\item
  report confidence intervals and failure rates;
\item
  compare naive n-step SAC with SACn-style corrections;
\item
  tune hyperparameters separately for each horizon;
\item
  test sequence-aware critics inspired by T-SAC.
\end{enumerate}

\section{Conclusion}\label{conclusion}

This project implements configurable n-step targets in SAC and tests
them in a small but reproducible ablation. Pendulum-v1 confirms that the
implementation works. HalfCheetah-v5 shows the main trade-off: \(n=3\)
helps, \(n=1\) remains a strong baseline, and \(n=5\) fails under the
same hyperparameters.

\begin{keybox}
\textbf{Final takeaway.} The n-step horizon is a stability parameter, not a bigger-is-better knob.
\end{keybox}

\section*{References}\label{references}
\addcontentsline{toc}{section}{References}

\protect\phantomsection\label{refs}
\begin{CSLReferences}{1}{0}
\bibitem[\citeproctext]{ref-espeholt2018impala}
Espeholt, Lasse, Hubert Soyer, Rémi Munos, Karen Simonyan, Volodymyr
Mnih, Tom Ward, Yotam Doron, et al. 2018. {``IMPALA: Scalable
Distributed Deep-RL with Importance Weighted Actor-Learner
Architectures.''} \emph{arXiv Preprint arXiv:1802.01561}.
\url{https://arxiv.org/abs/1802.01561}.

\bibitem[\citeproctext]{ref-gymnasium_halfcheetah}
Farama Foundation. 2026a. {``Gymnasium Documentation: HalfCheetah-V5.''}
\url{https://gymnasium.farama.org/environments/mujoco/half_cheetah/}.

\bibitem[\citeproctext]{ref-gymnasium_pendulum}
---------. 2026b. {``Gymnasium Documentation: Pendulum-V1.''}
\url{https://gymnasium.farama.org/environments/classic_control/pendulum/}.

\bibitem[\citeproctext]{ref-haarnoja2018sac}
Haarnoja, Tuomas, Aurick Zhou, Pieter Abbeel, and Sergey Levine. 2018.
{``Soft Actor-Critic: Off-Policy Maximum Entropy Deep Reinforcement
Learning with a Stochastic Actor.''} \emph{arXiv Preprint
arXiv:1801.01290}. \url{https://arxiv.org/abs/1801.01290}.

\bibitem[\citeproctext]{ref-haarnoja2018applications}
Haarnoja, Tuomas, Aurick Zhou, Kristian Hartikainen, George Tucker,
Sehoon Ha, Jie Tan, Vikash Kumar, et al. 2018. {``Soft Actor-Critic
Algorithms and Applications.''} \emph{arXiv Preprint arXiv:1812.05905}.
\url{https://arxiv.org/abs/1812.05905}.

\bibitem[\citeproctext]{ref-vanhasselt2018deadlytriad}
Hasselt, Hado van, Yotam Doron, Florian Strub, Matteo Hessel, Nicolas
Sonnerat, and Joseph Modayil. 2018. {``Deep Reinforcement Learning and
the Deadly Triad.''} \emph{arXiv Preprint arXiv:1812.02648}.
\url{https://arxiv.org/abs/1812.02648}.

\bibitem[\citeproctext]{ref-lyskawa2025sacn}
Łyskawa, Jakub, Jakub Lewandowski, and Paweł Wawrzyński. 2025. {``SACn:
Soft Actor-Critic with n-Step Returns.''} \emph{arXiv Preprint
arXiv:2512.13165}. \url{https://arxiv.org/abs/2512.13165}.

\bibitem[\citeproctext]{ref-munos2016retrace}
Munos, Rémi, Tom Stepleton, Anna Harutyunyan, and Marc G. Bellemare.
2016. {``Safe and Efficient Off-Policy Reinforcement Learning.''}
\emph{arXiv Preprint arXiv:1606.02647}.
\url{https://arxiv.org/abs/1606.02647}.

\bibitem[\citeproctext]{ref-sharma2017mixing}
Sharma, Sahil, Girish Raguvir J, Srivatsan Ramesh, and Balaraman
Ravindran. 2017. {``Learning to Mix n-Step Returns: Generalizing
Lambda-Returns for Deep Reinforcement Learning.''} \emph{arXiv Preprint
arXiv:1705.07445}. \url{https://arxiv.org/abs/1705.07445}.

\bibitem[\citeproctext]{ref-sutton2018rl}
Sutton, Richard S., and Andrew G. Barto. 2018. \emph{Reinforcement
Learning: An Introduction}. 2nd ed. MIT Press.

\bibitem[\citeproctext]{ref-tian2025chunking}
Tian, Dong, Onur Celik, and Gerhard Neumann. 2025. {``Chunking the
Critic: A Transformer-Based Soft Actor-Critic with n-Step Returns.''}
\emph{arXiv Preprint arXiv:2503.03660}.
\url{https://arxiv.org/abs/2503.03660}.

\end{CSLReferences}

\end{document}
